\documentclass[british,oneside,11pt]{amsart}
\usepackage[T1]{fontenc}
\usepackage{babel}
\usepackage{lmodern}
\usepackage[pdfusetitle,colorlinks]{hyperref}
\usepackage{bookml/bookml}
\bmlAltFormat{main.pdf}{PDF (serif)}
\hypersetup{pdfsubject={A compact BookML technical-book starter.}}

\title{A Data-Driven Technical Book}
\author{A. Researcher}
\date{2026}

\begin{document}
\maketitle
\tableofcontents

\section{Purpose}
BookML connects a LaTeX source tree to PDF and accessible web-oriented output
through LaTeXML. This starter keeps the document ordinary enough to compile as
a PDF while showing the small hooks that BookML adds for alternative formats.

\section{A source-first workflow}
Keep chapters and supporting data in separate files as the manuscript grows.
The BookML make rules can discover document sources, generate PDF variants,
and build HTML, EPUB, or SCORM outputs from the same source tree.

\subsection{A useful invariant}
Every generated format should preserve the same heading order, equations, and
cross-references. Treat the PDF as one checked projection of the source, not
as a separate manuscript.

\section{A worked example}
The accessible-book workflow can carry mathematical content as well:
\[
  f(x) = \int_{0}^{x} g(t)\,\mathrm{d}t,
  \qquad f'(x)=g(x).
\]
Add semantic descriptions to images with `\bmlDescription` and use the
BookML manual's HTML environments when a web-only enhancement is needed.

\section{Next steps}
Replace this text with chapters in the same directory and add a GNUmakefile
with `include bookml/bookml.mk` when you want BookML's multi-format build.

\iflatexml\else
  \bigskip
  \begin{center}
    {\small Copyright \copyright{} 2026 A. Researcher. Compiled with BookML.}
  \end{center}
\fi
\end{document}
